% !PRINTED 578
\weberpage{578}

\begin{center}
\S.\,158.

\gesperrt{Dedekind's Satz über die Körperdiscriminante.}
\end{center}

Die in der Discriminante $\Delta$ des K\"orpers $\Omega$ aufgehenden natürlichen Primzahlen haben in Bezug auf ihre Zerlegung in Primfactoren einen besonderen Charakter, über den ein Satz von \gesperrt{Dedekind} die bündigste Auskunft giebt, zu dessen Ableitung wir jetzt schreiten\footnote{\gesperrt{Dedekind}, \glqq Ueber die Discriminanten endlicher K\"orper\grqq{} im XXIX. Bande der Abhandlungen der Gesellschaft der Wissenschaften in G\"ottingen (1882).}.

Wir bilden nach \S.\,145, 2. die Potenzen der Basisform
\[
\tau = t_1\omega_1 + t_2\omega_2 + \cdots + t_n\omega_n,
\]
und erhalten die Ausdrücke:
\weq{1}{\tau^k = u_{1,k}\,\omega_1 + u_{2,k}\,\omega_2 + \cdots + u_{n,k}\,\omega_n,}
worin die $u_{i,k}$ ganze rationale Functionen der Variablen $t_1, t_2, \ldots, t_n$ sind. Wir setzen wie oben
\weq{2}{F(t) = N(t-\tau),}
so dass $F(\tau)=0$ ist. Nun bilden wir nach \S.\,144 die Discriminante
\weq{3}{\Delta(1,\tau,\tau^2,\ldots,\tau^{n-1}) = (-1)^{\frac{n(n-1)}{2}}\, N F'(\tau),}
wofür sich nach (1) der Werth $\Delta U^2$ ergiebt, wenn $\Delta$ die Körperdiscriminante und
\weq{4}{U = \left|\begin{array}{cccc} u_{1,0}, & u_{2,0}, & \cdots, & u_{n,0} \\ u_{1,1}, & u_{2,1}, & \cdots, & u_{n,1} \\ \cdots & \cdots & \cdots & \cdots \\ u_{1,n-1}, & u_{2,n-1}, & \cdots, & u_{n,n-1} \end{array}\right|,}
also eine ganze Function in $R$ ist.

Wir müssen nachweisen, dass die Form $U$ eine Einheit ist. Angenommen, es gehe in $U$ irgend eine Primzahl $p$ auf, so können wir, wie schon im \S.\,147 bewiesen ist, ein System ganzer Functionen $y_0, y_1, \ldots, y_{n-1}$ in $R$, \gesperrt{die nicht alle durch $p$ theilbar sind}, so bestimmen, dass für $i=1,2,\ldots,n$
\[
y_0\, u_{i,0} + y_1\, u_{i,1} + \cdots + y_{n-1}\, u_{i,n-1} \equiv 0\ (\mathrm{mod}\ p)
\]
wird. Dann aber ergiebt sich aus (1)
\[
y_0 + y_1\tau + \cdots + y_{n-1}\tau^{n-1} \equiv 0\ (\mathrm{mod}\ p).
\]
%CONT>>
